\documentclass[a4paper]{article}
%% Sets page size and margins
\usepackage[a4paper,top=2cm,bottom=2cm,left=2.8cm,right=2.8cm,marginparwidth=1.75cm]{geometry}

%% Useful packages
\usepackage{amsmath}
\usepackage{graphicx}
\usepackage[colorinlistoftodos]{todonotes}
\usepackage[colorlinks=true, allcolors=blue]{hyperref}

\title{Homework 2\\
600.482/682 Deep Learning\\
Fall 2025}
\date{}

\begin{document}
\maketitle
\vspace{-1.3cm}

\centerline{\textbf{Due Wednesday, Sep 17th, 2025, 11:59 pm EST}}
\vspace{0.2cm}

\noindent \textbf{Submission Instructions.} Please submit the following \textbf{two items} to Gradescope (entry code VWJRB3): 

\begin{enumerate}
    \item Your report (LaTeX generated PDF) named \texttt{homework2-report.pdf}.
    % 
    \item A Zip file containing: 
    % 
    (1) Your Python Jupyter Notebook (\texttt{homework2-notebook.ipynb}) and (2) an exported PDF of the notebook (with all cell outputs) to \texttt{homework2-notebook.pdf}.
    
    * On Google Colab, you may use ``File'' -\> ``Print'' -\> Save as PDF.
    
    * On Visual Studio, you may use ``Jupyter: Export to PDF'' in the Command palette.
\end{enumerate}


\vspace{5mm}

\noindent \textbf{Problems} 

\begin{enumerate}
% Problem 1
\item The goal of this problem is to minimize a function given a certain input using gradient descent by breaking down the overall function into smaller components via a computational graph. The function is defined as:
\begin{align*}
f(x_1, x_2, w_1, w_2) & = \sigma(w_1 x_1+w_2 x_2) + 0.5(w_1^2 + w_2^2) \\
& =  \frac{1}{1+e^{-(w_1 x_1+w_2 x_2)}} + 0.5(w_1^2 + w_2^2)\,.
\end{align*}
\begin{enumerate}
\item Please calculate $\frac{\partial f}{\partial x_1}, \frac{\partial f}{\partial x_2}, \frac{\partial f}{\partial w_1}, \frac{\partial f}{\partial w_2}$. 
% 
When computing the derivative for the sigmoid function $\sigma(x)$, please use sigmoid function $\sigma(x)$ in your expression.


\item Start with the following initialization: $w_1 = 0.2, w_2 = 0.4, x_1 = -0.4, x_2 = 0.5$, draw the computational graph. Please use backpropagation as we did in class. You may use the sigmoid function $\sigma(x)$ as a node in the graph.

You can draw the graph on paper and insert a photo into your report.

The goal is for you to practice working with computational graphs.

You must include the intermediate values in both forward and backward passes in your graph.

\item John would like to train a binary classification model that classifies whether a particular fruit is an apple or not based on two features: its color ($x_1$) and its roundness ($x_2$). 
% 
$w_1$ and $w_2$ are the weights he gives to the two features. 
% 
John tries to use the function $f(x_1, x_2, w_1, w_2)$ as the loss function for a single data point in his optimization problem. 
% 
Will this formulation of loss function work? 
% 
Briefly explain why or why not.
\end{enumerate}

% Problem 2
\item Doing the computations above is a lot of work, already for such a simple function as in Problem 1. 
% 
Clearly, it will be desirable to automate the forward and backward propagation process. 
% 
In contemporary frameworks, such as TensorFlow and PyTorch, backpropagation is handled automatically for all standard operations, which is very convenient. 
% 
This feature is commonly referred to as ``AutoGrad''. 
% 
It becomes possible due to one of the key observations we learned in class: difficult expressions can be broken down into simple sub-problems, the analytic gradients of which can be synthesized via chain rule. 
% 
In this problem, you will implement your own ``AutoGrad'' structure.
% 
\begin{enumerate}
\item 
Consider the Jupyter Notebook we provide (\textit{Homework2.ipynb}, which can be downloaded from \url{https://piazza.com/class_profile/get_resource/mdnhowckj2f5di/mf451vw4w4q1ci}) and follow the instructions therein to implement the missing operations and functionality. 
% 
To use Jupyter Notebook, we highly recommend uploading the \textit{Homework2.ipynb} file to Google Colab (\url{https://colab.research.google.com/}), where you may edit and run your notebook online. 
% 
If that option is not feasible for you, you can create a local Python 3.6 (or higher version) environment. 
% 
The only required external dependencies are \texttt{numpy}, \texttt{matplotlib}, and \texttt{jupyter} which you can install using the package manager of your choosing (\textit{e}.\textit{g}. \texttt{pip}, \texttt{conda}).
% 
\item Use the function and output derived in Problem 1 to test whether your ``AutoGrad'' structure works as intended. 


Attach a screenshot of your notebook's printed intermediate values for this test case (both during the forward and backward pass) to your report.

\end{enumerate}

	
% Problem 3: 2 dimensional data
\item 
Your goal in this exercise is to implement a linear classifier using the AutoGrad class you implemented in Problem 2, train it on a 2-dimensional toy dataset, and show the results tested on the {\bf same} dataset.
Recall that we formulate a linear classifier as $f(x, W, b) = Wx + b$. To quantify the model's performance, you will pass its output through a softmax function and then use Cross Entropy loss to measure agreement with the target output.\\

Write a Python program based on your AutoGrad structure that iteratively computes a small step in the direction of the negative gradient. The dataset (HW2\_Q3\_dataset) can be downloaded from \url{https://piazza.com/class_profile/get_resource/mdnhowckj2f5di/mf44x67wv1r47b}. 
{\bf Please go through the README.txt carefully before you start.}
Initialize your optimization with all elements of $W$ being small random numbers and $b=0$. As shown in the visualization of the data provided in the link above, you can reasonably expect to see your algorithm converge to a solution that linearly separates the given data samples.\\

\begin{enumerate}
	\item Please plot the following and attach them in the report:
	\begin{enumerate}
		\item Cross Entropy loss with respect to training iterations~(loss curve) 
		\item Training accuracy with respect to training iterations~(accuracy curve)
		\item Visualization of the ground truth labels and your model's decision boundaries~(decision boundary visualization)
	\end{enumerate}
	\item Report your final classification accuracy. Briefly discuss your result based on the visualization of data distribution.
\end{enumerate}

% Problem 4: 2D data convenient to be visualized
\item	In this problem, we would like you to compare the performance of single-layer and multi-layer classifiers on a different dataset, again using your AutoGrad structure.
	\begin{enumerate}
		\item Please use HW2\_Q4\_dataset (\url{https://piazza.com/class_profile/get_resource/mdnhowckj2f5di/mf44x33xayp450}) for this question. This data is two-dimensional. Create another linear model~(same as Problem 3) for this dataset and run your classification training. 
		
		Report your final classification accuracy. As in Problem 3, please attach plots of the i) loss curve, ii) accuracy curve, and iii) decision boundary visualization. Is this dataset linearly separable?
		
\item Now add one more layer of perceptions.
  
* Use ReLU as activation functions for the first hidden layer.

** Connect outputs with a softmax function. Initialize your optimization procedure with all elements of $W$ being small random numbers and $b=0$. Use Cross Entropy loss. Train this model on the same two-dimensional dataset in (a).
		
Report your classification accuracy. Attach plots of the i) loss curve, ii) accuracy curve, and iii) decision boundary visualization in the report.
		
\item Compare your results, including the classification accuracy and decision boundary plot, between 4(a) and 4(b). Briefly explain why introducing an additional layer causes these changes. 
	
\end{enumerate}

* If the output of a single layer is $f(x, W_1, b_1)$, then adding another layer basically means that the output is $g(f(x, W_1, b_1), W_2, b_2)$: the second layer uses the first layer's output as its input. Therefore, the output dimension of $f$ should match the expected input dimension of $g$. In our case, the output dimension from the ReLU activation in the first layer should match the input dimension of the second layer.
	
	
** The ReLU activation function is a nonlinear function defined as:
\[
  \textrm{ReLU}(x) = \textrm{max}(0, x)  
\]
If we use ReLU as our activation function for the first layer, instead of $f(x, W_1,b_1) = W_1x+b$, we have
\[
    f(x, W_1, b_1) = \textrm{ReLU}(W_1x+b) = \textrm{max}(0, W_1x+b)
\]
The non-linearity of the ReLU function adds non-linearity to the model.
    
\end{enumerate}

\end{document}















